% Nouveau EGA V (cote 25): archive cote 25, site page 163 (PDF page 164).
% typescript. Draft compared: none.
% AI-transcribed by Claude Opus 5.5 (claude-opus-5-5, adaptive thinking, effort high), Message Batches API, batch ega25b (msgbatch_01A7n5YuXubUPNK1zsepi8RG), 2026-11-13; original AI transcription, subsequently checked against both supplied scan parts; see PAGES_CHECKED.tsv and CHECK_DECISIONS.tsv.
\scanpage{25/163}
\mspage{\hand{6}}

un sous-préschéma X" de X, régulier, \var{\xout{tel que}}{}{typewriter strike-out}{ega25_25_163-v0001} contenant X, et tel que X
puisse se décrire dans X"\textsubscript{\var{\xout{\unsure{w}}}{}{typewriter strike-out}{ega25_25_163-v0002}} par une seule équation $\hand{\phi} = 0$ , admettant x comme
zéro \var{\inserted[typed]{sing}}{sing}{insertion}{ega25_25_163-v0003} élémentaire de multiplicité n.

\var{\inserted[typed]{On conclut en particulier que si A est local noeth élémentaire de}}{On conclut en particulier que si A est local noeth élémentaire de}{insertion}{ega25_25_163-v0004}
\var{\xout{15.7. On peut donner des variantes géométriques de 15.6\unsure{.,|,} \unsure{lorsqu'on}}}{}{typewriter strike-out}{ega25_25_163-v0005}
\var{\inserted[typed]{mult. n, alors A est Coh. Mac. \var{\xout{Il est clair}}{}{typewriter strike-out}{ega25_25_163-v0007} Nous montrerons au par.20}}{mult. n, alors A est Coh. Mac. \var{\xout{Il est clair}}{}{typewriter strike-out}{ega25_25_163-v0007} Nous montrerons au par.20}{insertion}{ega25_25_163-v0006}
\var{\xout{est sur un corps k (point de X, pas nécessairement rationnel sur k)}}{}{typewriter strike-out}{ega25_25_163-v0008}
(par un calcul facile d'éclatement) que le point fermé de Spec(A) est
le seul point singulier de Spec(A); il en résulte que \var{\xout{si dim}}{}{typewriter strike-out}{ega25_25_163-v0009} A est normal
sss dim A $\hand{\geqslant}$ 2, réduit sss dimA $\hand{\geqslant}$1. Si dim A = 0, alors  A est élémentaire
de multiplicité n sss  $m/m^2$ est de \var{\over{dim}{rang}}{rang}{writing over}{ega25_25_163-v0010}1 et  n est le plus petit entier
tel que $\underline{m}^n = 0$; on notera que ces anneaux sont aussi lesquotients d'anneaux
de valu\var{\xout{\unsure{at}}}{}{typewriter strike-out}{ega25_25_163-v0011}ion discrète par la\var{\inserted[typed]{\var{\xout{\unsure{ent}}}{}{typewriter strike-out}{ega25_25_163-v0013}}}{\var{\xout{\unsure{ent}}}{}{typewriter strike-out}{ega25_25_163-v0013}}{insertion}{ega25_25_163-v0012} puissance n.ème de \var{\deleted{l'}}{}{deletion}{ega25_25_163-v0014}\var{\inserted[typed]{\unsure{l|\hand{l}}eur}}{\unsure{l|\hand{l}}eur}{insertion}{ega25_25_163-v0015} idéal maximal.

15.7.   Il y a lieu de dégager les variantes ``géométriques'' de 15.6. ,
on trouvera en particulier que si X est un préschéma loc de type fini sur
le corps k, $x \hand{\in} X$, et si \var{\over{n}{d}}{d}{writing over}{ega25_25_163-v0016} $= \dim_x X$ , alors  x est une singularité \var{\xout{ordinaire}}{}{typewriter strike-out}{ega25_25_163-v0017}
\var{\inserted[typed]{élémentaire}}{élémentaire}{insertion}{ega25_25_163-v0018}
\var{\xout{\unsure{re}}}{}{typewriter strike-out}{ega25_25_163-v0019} de multiplicité  n de X, sss il existe un voisinage ouvert U de x
tel que U se plonge comme sous-préschéma d'un préschéma lisse X' sur k,
\var{\xout{de d}}{}{typewriter strike-out}{ega25_25_163-v0020} conne xe et de dimension \var{\over{n}{d}}{d}{writing over}{ega25_25_163-v0021}+1 , défini par une équation $\hand{\phi}$=0 , x étant
un zéro \var{\inserted[typed]{géom.}}{géom.}{insertion}{ega25_25_163-v0022} singulier ordinai\unsure{r|m}e de multiplicité n de $\hand{\phi}$ . On aura dit d'ailleurs
\var{\margin{$\surd$}}{\placedmargin{$\surd$}}{marginal note placed}{ega25_25_163-v0023} dans 15.5. que cela signifie  que  la \var{\inserted[typed]{valeur en x de}}{valeur en x de}{insertion}{ega25_25_163-v0024} partie principale $d^n\hand{\phi}$ de $\hand{\phi}$ , \var{\xout{restrei}}{}{typewriter strike-out}{ega25_25_163-v0025}
\var{\xout{nte à X=V($\hand{\phi}$)}}{}{typewriter strike-out}{ega25_25_163-v0026} (qui à priori est un\var{\xout{e section}}{}{typewriter strike-out}{ega25_25_163-v0027} \var{\inserted[typed]{élément}}{élément}{insertion}{ega25_25_163-v0028} de $\underline{P}^n_{X'/k}\hand{\otimes} \unsure{\underline{k}(x)}$ ,et plus
précisément de son Idéal d'augmentation $\underline{I}$ )  est en fait un\var{\xout{e section}}{}{typewriter strike-out}{ega25_25_163-v0029} \var{\inserted[typed]{élément de}}{élément de}{insertion}{ega25_25_163-v0030}
\var{\xout{de}}{}{typewriter strike-out}{ega25_25_163-v0031} $\underline{I}^n = \underline{Sym}^n_{\underline{O}_U}(\hand{\underline{\Omega}}^1_{X'/k}\hand{\otimes} \unsure{k(\hand{x})}\textsubscript{\var{\xout{\unsure{w}}}{}{typewriter strike-out}{ega25_25_163-v0032}})$), et  à ce titre est une \var{\xout{section}}{}{typewriter strike-out}{ega25_25_163-v0033} \var{\inserted[typed]{forme}}{forme}{insertion}{ega25_25_163-v0034} lisse \var{\xout{en x}}{}{typewriter strike-out}{ega25_25_163-v0035}
\var{\margin{|}}{\placedmargin{|}}{marginal note placed}{ega25_25_163-v0036} \hand{[}NB Oubli dans 15.1. : l'ensemble des points de lissité d'une forme est
\var{\margin{|}}{\placedmargin{|}}{marginal note placed}{ega25_25_163-v0037} ouvert\hand{]}. On remarquera de même qu'un tel point est isolé dans l'ensemble
des points non lisses, il est géom normal \var{\xout{\illegible{1}}}{}{typewriter strike-out}{ega25_25_163-v0038} (resp. géom réduit) sss
\var{\over{n}{d}}{d}{writing over}{ega25_25_163-v0039} $\hand{\geqslant} 2$  (resp. \var{\over{n}{d}}{d}{writing over}{ega25_25_163-v0040} $\hand{\geqslant} 1$). Si \var{\over{n}{d}}{d}{writing over}{ega25_25_163-v0041} = 0 et k(x) = k (\var{\inserted[typed]{i.e. x  isolé et}}{i.e. x  isolé et}{insertion}{ega25_25_163-v0042} x rationnel sur k) alors \var{\xout{x}}{}{typewriter strike-out}{ega25_25_163-v0043}
point singulier élémentaire  de mult\hand{\unsure{,}} \unsure{n} signifie que $\underline{O}_{X,x}$  est  \var{\deleted{k-isomor}}{}{deletion}{ega25_25_163-v0044}

\mspage{\hand{163}}
